% M1031
\section{Результати числового аналізу}
% M1032
\label{sec:numerical-results}

% M1034
Числовий аналіз проводимо для нормованих параметрів $\,d_i/l$, $\delta_i'$ й $J_i'$, визначених у розділі~\ref{sec:derivation}. Фізично допустимими вважаємо лише ті комбінації матеріалу, кута та знака навантаження, для яких обидві нерівності
% M1038
\begin{equation*}
  Cg_2<0,\qquad CQ_i>0
\end{equation*}
% M1041
виконуються одночасно. Перша нерівність забезпечує стискальний контакт між берегами V-подібної міжфазної тріщини зсуву, а друга — розтягувальне нормальне напруження вздовж бісектриси кутового сектора матеріалу, в якому утворюється зона передруйнування.

% M1047
Щоб з’ясувати, чи допустимість зони передруйнування лише в одному з матеріалів, виявлена для базової пари, властива саме цьому прикладу, додатково провели широкий числовий аналіз знаків. Оскільки у виразі для $q_i$ у додатку~\ref{app:asymptotics} множник для $i=1,2$ є спільним, а відрізняється лише функція в чисельнику $S_i$, маємо
% M1052
\begin{equation*}
  \operatorname{sgn}(Q_1Q_2)
  =\operatorname{sgn}(S_1S_2).
\end{equation*}
% M1056
Отже, одночасна допустимість за фіксованого знака $C$ вимагала б $S_1S_2>0$. Детермінований перебір 92~800 невироджених комбінацій охопив $10^{-3}\le E_1/E_2\le10^3$, по вісім незалежно варійованих значень коефіцієнта Пуассона між $-0.95$ та $0.49$ для кожного матеріалу і $\alpha=3^\circ,6^\circ,\ldots,177^\circ$, за винятком границі прямолінійної межі поділу. Випадків із $S_1S_2>0$ не виявлено. Окреме дослідження випадкової вибірки з 20~000 додаткових точок із фіксованим початковим станом генератора, де $E_1/E_2$ вибирали рівномірно в логарифмічному масштабі на $10^{-3}\le E_1/E_2\le10^3$, а також $-0.99<\nu_i<0.499$ та $0.2^\circ<\alpha<179.8^\circ$, так само не виявило жодного такого випадку. У двох точках вибірки, найближчих до $S_1S_2=0$, від’ємний знак зберігся після повторного обчислення з 50 десятковими цифрами точності. Отже, у цій широкій області параметрів за заданого знака навантаження умову існування розтягувальної зони задовольняє щонайбільше один матеріал. Ці числові дані не є доведенням глобальної взаємовиключності поза дослідженою областю.

% M1072
Для базової пари матеріалів
% M1073
\begin{equation*}
  E_1/E_2=0.5,\qquad \nu_1=\nu_2=0.3,
\end{equation*}
% M1076
зміни знака коефіцієнта $g_2$ відбуваються при
% M1077
\begin{equation*}
  \alpha_1=12.92^\circ,\qquad
  \alpha_2=107.09^\circ.
\end{equation*}
% M1081
Ці значення отримано безпосереднім розв’язанням $D_0(-1-\lambda_0)=0$ разом з $g_2=0$. За $C>0$ зона передруйнування фізично допустима в матеріалі~2 при $0<\alpha<\alpha_1$ і в матеріалі~1 при $\alpha_2<\alpha<180^\circ$. За $C<0$ відповідні інтервали становлять $\alpha_1<\alpha<90^\circ$ для матеріалу~1 і $90^\circ<\alpha<\alpha_2$ для матеріалу~2. Значення $\alpha=90^\circ$ відповідає прямолінійній межі поділу й далі розглядається як вироджена аналітична границя, без трактування як звичайної внутрішньої точки числової гілки.

% M1091
Щоб оцінити вплив пружного контрасту, ті самі два перехідні кути обчислено для діапазону $E_1/E_2<1$ за $\nu_1=\nu_2=0.3$. Результати наведено в табл.~\ref{tab:g2-transition-angles}. Зі збільшенням $E_1/E_2$ перший перехідний кут $\alpha_1$ монотонно зростає, а $\alpha_2$ монотонно спадає. Отже, пружний контраст систематично зміщує межі кутових інтервалів фізичної допустимості.

% M1099
\begin{table}[htbp]
% M1100
\centering
% M1101
\caption{Кути, за яких коефіцієнт $g_2$ змінює знак, при $\nu_1=\nu_2=0.3$.}
% M1103
\label{tab:g2-transition-angles}
% M1104
\scriptsize
% M1105
\setlength{\tabcolsep}{2.6pt}
% M1106
\begin{tabular}{@{}lrrrrrrrrrrr@{}}
% M1107
\toprule
% M1108
$E_1/E_2$ & 0.01 & 0.10 & 0.20 & 0.30 & 0.40 & 0.50 & 0.60 & 0.70 & 0.80 & 0.90 & 0.99 \\
% M1109
\midrule
% M1110
$\alpha_1$, $^\circ$ & 2.58 & 7.03 & 9.25 & 10.77 & 11.95 & 12.92 & 13.74 & 14.45 & 15.08 & 15.65 & 16.11 \\
% M1111
$\alpha_2$, $^\circ$ & 108.02 & 107.90 & 107.72 & 107.52 & 107.30 & 107.09 & 106.88 & 106.69 & 106.50 & 106.32 & 106.17 \\
% M1112
\bottomrule
% M1113
\end{tabular}
% M1114
\end{table}

% M1117
Випадки з $E_1/E_2>1$ взаємопов’язані. Їх одержуємо взаємною заміною матеріалів за правилом
% M1119
\begin{equation*}
  (E_1,\nu_1)\leftrightarrow(E_2,\nu_2),
  \qquad
  \alpha\leftrightarrow\pi-\alpha,
  \qquad
  1\leftrightarrow2.
\end{equation*}
% M1126
Для відповідних локальних величин це перетворення також дає $d_1\leftrightarrow d_2$ і $\lambda_1\leftrightarrow\lambda_2$, а для перехідних півкутів — $\alpha_1\leftrightarrow\pi-\alpha_2$. Для $\nu_1=\nu_2$ це охоплює заміну $E_1/E_2\leftrightarrow E_2/E_1$, тому в табл.~\ref{tab:g2-transition-angles} достатньо навести $E_1/E_2<1$.

% M1134
\subsection{Залежність від нормованого навантаження}
% M1135
\label{subsec:figure2}

% M1137
На рис.~\ref{fig:load-dependence} показано параметри зони передруйнування як функції
% M1139
\begin{equation*}
  \sigma'=\frac{Cl^{\lambda_0}}{\sigma_i}
\end{equation*}
% M1142
для чотирьох базових конфігурацій: $\alpha=10^\circ$ ($i=2$, $C>0$), $45^\circ$ ($i=1$, $C<0$), $105^\circ$ ($i=2$, $C<0$) та $135^\circ$ ($i=1$, $C>0$). Для кожної конфігурації показано лише фізично допустиму піввісь параметра навантаження.

% M1149
\resultfigure{Fig2.pdf}{Залежність нормованих параметрів зони передруйнування від безрозмірного параметра навантаження $\sigma'$ для чотирьох базових кутів і фізично допустимих матеріалів.}{fig:load-dependence}

% M1154
У всіх чотирьох випадках $d_i/l$, $\delta_i'$ й $J_i'$ монотонно зростають зі збільшенням $|\sigma'|$. Ця спільна тенденція безпосередньо випливає з аналітичних співвідношень $\delta_i'\propto d_i/l$ та $J_i'\propto d_i/l$ за фіксованої геометрії. Серед чотирьох наведених базових конфігурацій найбільші значення всіх трьох параметрів за $|\sigma'|=0.5$ отримано для $\alpha=45^\circ$, $C<0$, а саме
% M1161
\begin{equation*}
  \frac{d_1}{l}=0.1474,\qquad
  \delta_1'=0.1458,\qquad
  J_1'=0.0836.
\end{equation*}
% M1166
Ця тенденція також стосується маломасштабного припущення. Зі збільшенням навантаження зростають локальні деформаційний та енергетичний параметри, а модель наближається до межі своєї асимптотичної застосовності.

% M1170
За цими кривими навантаження можна безпосередньо визначити критичне навантаження ініціювання відгалуження, якщо критичний параметр матеріалу задано незалежно. Якщо матеріал характеризується критичним розкриттям $\delta_{i\mathrm{c}}$ або критичним зональним енергетичним параметром $J_{i\mathrm{c}}$, перетин із відповідною кривою визначає нормоване критичне навантаження $\sigma'_c$ і, отже, амплітуду зовнішнього поля,
% M1177
\begin{equation*}
  C_c=\sigma'_c\sigma_i l^{-\lambda_0}.
\end{equation*}
% M1180
У наведених прикладах числових значень $\delta_{i\mathrm{c}}$ або $J_{i\mathrm{c}}$ не задано, тому графіки не визначають конкретної точки ініціювання. Формула~\eqref{eq:energy-opening-relation} показує, що за фіксованої геометрії та властивостей матеріалів $\delta_i$ і $J_i$ є лінійно пов’язаними мірами тієї самої рівноважної довжини зони. Отже, за $J_{i\mathrm{c}}=\chi_i\sigma_i\delta_{i\mathrm{c}}$ деформаційний і енергетичний критерії визначають той самий критичний стан.

% M1189
\subsection{Залежність від півкута зламу межі поділу}
% M1190
\label{subsec:figure3}

% M1192
На рис.~\ref{fig:angle-dependence} показано залежність нормованих параметрів зони передруйнування від півкута $\alpha$ для тієї самої базової пари матеріалів при $\sigma'=\pm0.5$. Залишено лише ті поєднання матеріалу й знака навантаження, які задовольняють умови фізичної допустимості. Переходи між допустимими гілками відбуваються при
% M1197
$\alpha_1$, $90^\circ$ і $\alpha_2$.

% M1199
\resultfigure{Fig3.pdf}{Залежність нормованих параметрів зони передруйнування від півкута $\alpha$ при $\sigma'=\pm0.5$. Показано лише фізично допустимі поєднання матеріалу та знака навантаження. Значення $\alpha=90^\circ$ є аналітичною границею.}{fig:angle-dependence}

% M1206
На двох широких допустимих гілках для матеріалу~1 кутові залежності мають виражені внутрішні максимуми, а на коротких гілках для матеріалу~2 найбільші значення досягаються поблизу меж допустимості. Максимуми $d_i/l$, $\delta_i'$ й $J_i'$ близькі, але відповідають різним кутам. Для гілки матеріалу~1 з $C<0$ перебір із кроком один градус дає
% M1211
\begin{equation*}
\begin{aligned}
  \max(d_1/l)&=0.1779 &&\text{при }\alpha=35^\circ,\\
  \max\delta_1'&=0.2077 &&\text{при }\alpha=32^\circ,\\
  \max J_1'&=0.1139 &&\text{при }\alpha=33^\circ.
\end{aligned}
\end{equation*}
% M1218
Близькість максимумів узгоджується з тим, що за кожного фіксованого $\alpha$ величини $\delta_i'$ і $J_i'$ пропорційні $d_i/l$, а самі коефіцієнти пропорційності залежать від кута. Отже, рис.~\ref{fig:angle-dependence} кількісно показує вплив геометрії на деформаційну та енергетичну характеристики за заданого навантаження. Щоб з цих кривих отримати кутову залежність критичного прикладеного навантаження, потрібно задати відповідний критичний параметр матеріалу й розв’язати рівняння критерію окремо для кожного $\alpha$.

% M1227
Максимум гілки матеріалу~1 з $C<0$ на рис.~\ref{fig:angle-dependence} відповідає $d_i/l$ порядку $0.18$. Це значення вже не можна однозначно вважати малим порівняно з одиницею, тоді як локальна постановка спирається на припущення $d_i/l\ll1$. Тому $d_i/l<0.1$ приймаємо за основну область кількісної інтерпретації, в якій маломасштабне наближення найкраще обґрунтоване. Ділянки кривих із $d_i/l>0.1$ збережено, щоб показати структуру й тенденції аналітичного розв’язку. Їх слід тлумачити як екстраполяції локальної моделі з нижчою кількісною точністю.

% M1238
Числове обчислення плюс-факторів для рис.~\ref{fig:angle-dependence} потребує особливої уваги при $\alpha\to0$ і $\alpha\to180^\circ$. Поблизу цих границь за фіксованого контуру факторизації збіжність повільна. Тому в розрахунках контур адаптивно подовжували. Подальше подовження хвоста змінило обчислені плюс-фактори щонайбільше на $8.69\times10^{-13}$. Отже, рис.~\ref{fig:angle-dependence} побудовано за збіжними значеннями, отриманими з адаптивним подовженням.

% M1247
\subsection{Показники сингулярності та фізична допустимість}
% M1248
\label{subsec:figure4}

% M1250
Для порівняння структури локальних полів напружень розглянемо три математичні конфігурації тієї самої околиці кутової точки: кутову точку непошкодженої межі поділу двох матеріалів, кутову точку із заданою стаціонарною міжфазною тріщиною зсуву та конфігурацію, що містить і цю тріщину, і локальну зону передруйнування. Для базової пари матеріалів відповідні порядки сингулярності описують показники $\lambda_0$, $\lambda$ й $\lambda_i$. Порівняння характеризує відмінності моделей локального поля без припущення про хронологічну послідовність розвитку руйнування.

% M1259
\resultfigure{Fig4.pdf}{Показники сингулярності та фізично допустимі ділянки для базової пари матеріалів. Повні математичні гілки й фізично допустимі ділянки розрізнено за умовами $Cg_2<0$ та $CQ_i>0$.}{fig:singularity-exponents}

% M1266
Для базової пари матеріалів на фізично допустимих ділянках сингулярність конфігурації з міжфазною тріщиною зсуву істотно сильніша, ніж у непошкодженій кутовій точці: $|\lambda|$ перевищує $|\lambda_0|$. Введення зони передруйнування частково послаблює цю сингулярність, і для відповідного матеріалу
% M1271
\begin{equation*}
  |\lambda_0|<|\lambda_i|<|\lambda|.
\end{equation*}
% M1274
Отже, зона передруйнування дає локальний ефект екранування: знижує порядок сингулярності, пов’язаної з міжфазною тріщиною зсуву, але не усуває концентрації напружень повністю. Крім того, на всіх відкритих інтервалах фізичної допустимості базового розрахунку
% M1278
\begin{equation*}
  -\frac12<\lambda_i<0,
\end{equation*}
% M1281
тому залишкова сингулярність у $O$ слабша за класичну сингулярність $r^{-1/2}$ для тріщини. Обидва результати встановлено для розглянутої пари матеріалів та її фізично допустимих гілок. Для довільного пружного контрасту потрібна додаткова перевірка.

% M1287
Сильніша сингулярність конфігурації з міжфазною тріщиною створює сприятливі умови для локального руйнування в одному з матеріалів. Водночас існування та розмір зони передруйнування також залежать від умов $Cg_2<0$ й $CQ_i>0$, опору відриву $\sigma_i$ і рівня зовнішнього навантаження. Отже, показники $\lambda_0$, $\lambda$ і $\lambda_i$ характеризують підсилення й екранування локального поля, а ініціювання відгалуження пов’язане з досягненням критичного $\delta_i$ або $J_i$ уздовж відповідної сім’ї рівноважних станів.

% M1298
\section{Обговорення}
% M1299
\label{sec:discussion}

% M1301
\subsection{Сім’я квазістатичних станів та ініціювання відгалуження}
% M1302
\label{subsec:discussion-quasistatic}

% M1304
У цій постановці V-подібна міжфазна тріщина зсуву є заданою стаціонарною конфігурацією: довжина $l$ кожної гілки тріщини залишається сталою, а рівноважна довжина зони передруйнування $d_i$ змінюється з параметром зовнішнього навантаження $C$. Отже, розв’язок визначає однопараметричну сім’ю квазістатичних станів
% M1309
\begin{equation*}
  C\longrightarrow d_i(C)
  \longrightarrow
  \bigl\{\delta_i(C),W_i(C),J_i(C)\bigr\}.
\end{equation*}
% M1314
Термін «квазістатична» тут стосується сім’ї рівноважних конфігурацій зони, у яких нехтуємо інерцією.

% M1317
Цю сім’ю можна безпосередньо використати для визначення критичного навантаження ініціювання відгалуження. Якщо для матеріалу незалежно задано критичне розкриття $\delta_{i\mathrm{c}}$ або критичний зональний енергетичний параметр $J_{i\mathrm{c}}$, критичний стан визначається умовою
% M1321
\begin{equation*}
  \delta_i(C_c)=\delta_{i\mathrm{c}}
  \qquad\text{або}\qquad
  J_i(C_c)=J_{i\mathrm{c}}.
\end{equation*}
% M1326
Коли відома відповідна критична характеристика матеріалу, ці співвідношення визначають $C_c$. Подальше поширення й скінченна траєкторія відгалуженої тріщини потребують окремої постановки.

% M1330
\subsection{Конкурентні механізми руйнування}
% M1331
\label{subsec:discussion-competing}

% M1333
Для стаціонарної V-подібної тріщини можливі два різні локальні механізми. Перший пов’язаний із її кутовою точкою $O$. Сингулярне поле кутової точки підтримує розтягувальну зону передруйнування в одному з матеріалів, а досягнення критичного $\delta_i$ або $J_i$ означає ініціювання відгалуження. Другий механізм локалізований у вершинах гілок V-подібної тріщини. Для цих вершин Назаренко й Кіпніс~\citeyearpar{NazarenkoKipnis2022Equilibrium} отримали коефіцієнт інтенсивності напружень за модою~II, який у наших позначеннях має вигляд
% M1341
\begin{equation}
\label{eq:branch-tip-KII}
  K_{\mathrm{II}}
  =g_1 k\,C\,l^{\lambda_0+1/2},
\end{equation}
% M1346
де
% M1347
\begin{equation*}
  k=
  \frac{
    2\sqrt{2}(1+\kappa_2e)\Gamma(\lambda_0+3/2)
  }{
    [1+\kappa_1+(1+\kappa_2)e]\,
    \Gamma(\lambda_0+2)G^+(-1-\lambda_0)
  }.
\end{equation*}
% M1356
Тут $G^+$ — плюс-фактор базового ядра $G$, визначеного в додатку~\ref{app:asymptotics}. Їхня умова міжфазного руйнування має вигляд $K_{\mathrm{II}}=K_{\mathrm{IIc}}$. Для абсолютного значення навантаження це дає
% M1359
\begin{equation*}
  |C_{\mathrm{tip}}|
  =
  \frac{K_{\mathrm{IIc}}}
  {|g_1k|\,l^{\lambda_0+1/2}}.
\end{equation*}

% M1366
Для відгалуження з кутової точки нормоване критичне навантаження $\sigma'_c$ визначаємо або з $\delta_i=\delta_{i\mathrm{c}}$, або з $J_i=J_{i\mathrm{c}}$ за формулами~\eqref{eq:normalized-process-zone-length}--\eqref{eq:process-zone-energy-parameter}. Оскільки $\sigma'=Cl^{\lambda_0}/\sigma_i$, відповідне абсолютне значення навантаження дорівнює
% M1373
\begin{equation*}
  |C_{\mathrm{corner}}|
  =
  \sigma_i l^{-\lambda_0}|\sigma'_c|.
\end{equation*}
% M1378
Ввівши безрозмірну міжфазну тріщиностійкість за модою~II
% M1379
\begin{equation*}
  \mathcal R_{\mathrm{II}}
  =
  \frac{K_{\mathrm{IIc}}}{\sigma_i\sqrt{l}},
\end{equation*}
% M1384
одержуємо співвідношення між двома критичними навантаженнями
% M1385
\begin{equation}
\label{eq:mechanism-competition-ratio}
  \frac{|C_{\mathrm{corner}}|}{|C_{\mathrm{tip}}|}
  =
  \frac{|g_1k|\,|\sigma'_c|}{\mathcal R_{\mathrm{II}}}.
\end{equation}
% M1391
Отже, порога відгалуження з кутової точки досягаємо першим за умови
% M1392
\begin{equation*}
  |g_1k|\,|\sigma'_c|<\mathcal R_{\mathrm{II}},
\end{equation*}
% M1395
а за протилежної нерівності першим досягаємо порога поширення від вершин гілок уздовж межі поділу. Рівність визначає межу між областями реалізації конкурентних механізмів.

% M1399
Для базового варіанта $E_1/E_2=0.5$, $\nu_1=\nu_2=0.3$ та $\alpha=45^\circ$ незалежно обчислений коефіцієнт дорівнює $g_1k=-0.6794$. Якщо критерій відгалуження з кутової точки виразити через нормоване критичне розкриття
% M1404
\begin{equation*}
  \Delta_1=
  \frac{E_1\delta_{1\mathrm{c}}}
  {4(1-\nu_1^2)\sigma_1l},
\end{equation*}
% M1409
межа між областями реалізації конкурентних механізмів набуває вигляду
% M1410
\begin{equation*}
  \mathcal R_{\mathrm{II}}^*
  =1.1064\,\Delta_1^{0.6132}.
\end{equation*}
% M1414
Рівнозначно, за $\mathcal J_1=E_1J_{1\mathrm{c}}/
[4(1-\nu_1^2)\sigma_1^2l]$,
% M1417
\begin{equation*}
  \mathcal R_{\mathrm{II}}^*
  =1.5558\,\mathcal J_1^{0.6132}.
\end{equation*}
% M1421
Обидва подання визначають ту саму межу за $J_{1\mathrm{c}}=\chi_1\sigma_1\delta_{1\mathrm{c}}$.

% M1424
Отже, відгалуження з кутової точки та поширення від вершин гілок уздовж межі поділу можна кількісно порівняти як конкурентні механізми за тієї самої стаціонарної геометрії, без заданої послідовності стадій. Для однозначного числового визначення порядку їх реалізації в конкретній системі матеріалів потрібні незалежно виміряні характеристики опору руйнуванню $K_{\mathrm{IIc}}$ і або $\delta_{i\mathrm{c}}$, або $J_{i\mathrm{c}}$.

% M1431
Маломасштабне припущення $d_i\ll l$ дає змогу в провідному наближенні розглядати обидві локальні області окремо. Поле стаціонарної V-подібної тріщини є зовнішнім полем для задачі про зону передруйнування поблизу $O$, а зворотним впливом малої зони на поля у вершинах гілок нехтуємо. Коли $d_i/l$ перестає бути малим, потрібна повністю зв’язана крайова задача.

% M1438
\subsection{Енергетичний і деформаційний критерії}
% M1439
\label{subsec:discussion-criteria}

% M1441
Згідно з формулою~\eqref{eq:energy-opening-relation}, за фіксованих геометрії та пружних властивостей
% M1443
\begin{equation*}
  J_i=\chi_i\sigma_i\delta_i.
\end{equation*}
% M1446
Отже, $J_i$ і $\delta_i$ є залежними характеристиками поточного стану зони, оскільки обидві однозначно визначаються тією самою рівноважною довжиною $d_i$. Тому деформаційний критерій $\delta_i=\delta_{i\mathrm{c}}$ та енергетичний критерій $J_i=J_{i\mathrm{c}}$ визначають той самий критичний стан, якщо
% M1451
\begin{equation*}
  J_{i\mathrm{c}}
  =\chi_i\sigma_i\delta_{i\mathrm{c}}.
\end{equation*}

% M1456
Ця еквівалентність справджується для розглянутої однопараметричної моделі за фіксованої геометрії. Вона не означає універсальної тотожності між двома незалежно визначеними критичними характеристиками в довільних задачах руйнування, оскільки $\chi_i$ залежить від геометрії кутової точки та пружного контрасту. У застосуваннях можна використовувати будь-яку з двох критичних характеристик за умови, що її встановлено незалежно, а два критерії не накладають як незалежні обмеження на той самий стан моделі.

% M1464
\subsection{Залишкове поле кутової точки й локальний потік енергії}
% M1465
\label{subsec:discussion-flux}

% M1467
У конфігурації із зоною передруйнування поле в кутовій точці $O$ залишається сингулярним. Якщо $H_i$ — амплітуда провідного власного поля кутової точки, то
% M1470
\begin{equation*}
  \sigma\sim H_i r^{\lambda_i}.
\end{equation*}
% M1473
Відповідна густина енергії пружної деформації має порядок
% M1474
\begin{equation*}
  w\sim\frac{H_i^2}{E_i}r^{2\lambda_i},
\end{equation*}
% M1477
а локальний конфігураційний потік енергії через контур характерного радіуса $R$ змінюється за законом
% M1479
\begin{equation}
\label{eq:local-elastic-flux-scaling}
  J_i^{(\mathrm{el})}(R)
  \sim\frac{H_i^2}{E_i}R^{2\lambda_i+1}.
\end{equation}

% M1485
У фізично допустимих режимах, установлених числовим аналізом, $-1/2<\lambda_i<0$. Отже, $2\lambda_i+1>0$ і
% M1487
\begin{equation*}
  \lim_{R\to0}J_i^{(\mathrm{el})}(R)=0.
\end{equation*}
% M1490
Тому залишкова сингулярність напружень у $O$ не створює скінченного ненульового локального пружного внеску в енергетичний параметр зони. У цьому сенсі сингулярність кутової точки в конфігурації із зоною передруйнування слабша за класичну сингулярність $r^{-1/2}$ для тріщини.

% M1495
Ця відмінність також визначає межу застосовності адитивного представлення, використаного Камінським і Дудиком~\citeyearpar{KaminskyDudyk2026Criteria}. Для розглянутої в тій роботі наявної тріщини дійсна частина показника сингулярності зовнішнього поля дорівнює $-1/2$, і скінченний потік пружної енергії може входити до загального балансу енергії разом із внеском зони передруйнування. У цій задачі для кутової точки з $\lambda_i>-1/2$ випливає $J_i^{(\mathrm{el})}(R)\to0$, тому в критерії ініціювання відгалуження окремий скінченний пружний доданок $J_i^{(\mathrm{el})}$ не потрібен. Отже, зональний параметр $J_i=\dd W_i/\dd d_i$ можна застосовувати без такої адитивної поправки. Це не надає $J_i$ змісту незалежно встановленого глобального контурного інтеграла.

% M1507
\subsection{Область застосовності та узагальнення}
% M1508
\label{subsec:discussion-scope}

% M1510
Аналіз ґрунтується на розділенні масштабів $d_i\ll l\ll L$. Основною областю кількісного застосування вважаємо $d_i/l<0.1$. Ділянки числових кривих, де $d_i/l$ перевищує $0.1$ і досягає приблизно $0.18$ поблизу максимумів, корисні для демонстрації тенденцій і є екстраполяціями локального наближення. У цих режимах зростає роль членів вищих порядків зовнішнього асимптотичного поля та взаємодії локальних областей поблизу кутової точки й вершин гілок тріщини.

% M1519
Локальний характер розв’язку також полегшує його узагальнення на ширший клас задач. Для розвиненої симетричної V-подібної міжфазної тріщини зі скінченними гілками множник $CQ_i$, що задає амплітуду зовнішнього поля в локальній задачі, можна замінити амплітудою, або узагальненим коефіцієнтом інтенсивності напружень, відповідного сингулярного поля кутової точки. Таке узагальнення узгоджене за умови, що зона передруйнування мала й застосовна та сама локальна задача на власні значення. Тому модель можна використовувати для аналізу ініціювання відгалуження в точці зламу межі поділу. Споріднене застосування — розколювання тіла гострим клином, коли його локальна крайова задача зводиться до аналогічної кутової конфігурації.

% M1531
\section{Висновки}
% M1532
\label{sec:conclusions}

% M1534
\begin{enumerate}
% M1535
\item Побудовано аналітичну модель ініціювання відгалуження в кутовій точці V-подібної міжфазної тріщини зсуву. Сингулярне поле заданої стаціонарної тріщини використано як зовнішнє поле для малої розтягувальної зони передруйнування вздовж бісектриси одного з матеріалів. Одержано замкнені співвідношення для рівноважної довжини зони $d_i$, розкриття $\delta_i$, роботи $W_i$ сил опору відриву та зонального енергетичного параметра $J_i=\dd W_i/\dd d_i$.

% M1544
\item Фізична допустимість локальної зони передруйнування вимагає одночасно стискального контакту берегів V-подібної міжфазної тріщини, $Cg_2<0$, і розтягувального поля нормальних напружень у матеріалі, що містить зону, $CQ_i>0$. Під час широкого числового аналізу 92~800 детермінованих комбінацій параметрів і 20~000 додаткових випадково вибраних комбінацій не знайдено конфігурації, в якій за того самого знака навантаження обидва матеріали задовольняли б умову розтягувальної зони. Отже, допустимість зони передруйнування лише в одному з матеріалів підтверджено чисельно в зазначеній області, без твердження про загальну теорему єдиності. Кут межі поділу та пружний контраст зміщують межі допустимих інтервалів і змінюють рівноважні характеристики зони.

% M1556
\item Зона передруйнування послаблює сингулярність кутової точки конфігурації, хоча й не усуває її. На фізично допустимих інтервалах для базової пари матеріалів $-1/2<\lambda_i<0$. Отже, залишкове поле слабше за класичну сингулярність $r^{-1/2}$ для тріщини, а його локальний конфігураційний потік пружної енергії прямує до нуля під час стягування контуру до кутової точки. Тому в локальному критерії ініціювання відгалуження для цієї кутової задачі зональний енергетичний параметр $J_i$ не потребує окремого скінченного пружного доданка.

% M1566
\item За фіксованої геометрії розкриття зони та зональний енергетичний параметр однозначно пов’язані через ту саму рівноважну довжину зони: $J_i=\chi_i\sigma_i\delta_i$. Отже, деформаційний критерій $\delta_i=\delta_{i\mathrm{c}}$ і енергетичний критерій $J_i=J_{i\mathrm{c}}$ визначають той самий критичний стан лише за узгодженого калібрування незалежно заданих критичних параметрів, $J_{i\mathrm{c}}=\chi_i\sigma_i\delta_{i\mathrm{c}}$. Якщо будь-який із цих критичних параметрів матеріалу задано незалежно, модель визначає відповідне критичне зовнішнє навантаження ініціювання відгалуження.

% M1578
\item Кількісне застосування моделі спирається на маломасштабне припущення $d_i\ll l$. За основну область застосовності прийнято $d_i/l<0.1$, а ділянки кривих, що сягають приблизно $d_i/l=0.18$, слід тлумачити як екстраполяції локального наближення. Відгалуження з кутової точки й поширення від вершин гілок V-подібної тріщини вздовж межі поділу є конкурентними локальними механізмами, критичні навантаження яких задовольняють формулу~\eqref{eq:mechanism-competition-ratio}. Отже, порядок їх реалізації визначається співвідношенням нормованої міжфазної тріщиностійкості за модою~II $K_{\mathrm{IIc}}/(\sigma_i\sqrt l)$ і порога відгалуження з кутової точки $|g_1k|\,|\sigma'_c|$. Локальність постановки дає змогу узагальнити її на розвинені V-подібні міжфазні тріщини, замінивши амплітуду зовнішнього поля амплітудою відповідного поля кутової точки. Споріднене застосування — розколювання тіла гострим клином.
% M1593
\end{enumerate}
